%----------------------------------------------------------------------
\section{Detailed Verifications}\label{app:b}
%----------------------------------------------------------------------

This appendix gives component-by-component
verifications for the non-trivial propositions of
the main text. The level of detail is intended to
suffice for mechanization in a proof assistant.

%----------------------------------------------------------------------
\subsection{XOR-bridge correctness}\label{app:b-xor-bridge}
%----------------------------------------------------------------------

Full verification of Proposition \ref{prop:xor-bridge}.

\emph{Setup.} Let $L : \GF(2)^N \to \GF(2)^N$ be
$\GF(2)$-linear, $m \in \GF(2)^N$ a mask, $r_0
\in \GF(2)^N$ the initial register state.

\emph{Goal.} After the three-step micro-kernel of
Definition \ref{def:xor-bridge}, show
$r = L(r_0 \wedge m) \oplus (r_0 \wedge \neg m)$.

\emph{Step 1.} $w \leftarrow r_0 \wedge \neg m$.
Bitwise: $w_i = r_{0,i}$ if $m_i = 0$, else $w_i = 0$.

\emph{Step 2.} $r \leftarrow L(r_0)$. The register
holds $L$ applied componentwise (in the sense of $L$'s
linear structure) to the full $r_0$.

\emph{Step 3.} $r \leftarrow r \oplus L(w) \oplus w$.
Compute:
\begin{align*}
  r &= L(r_0) \oplus L(w) \oplus w \\
    &= L(r_0 \oplus w) \oplus w
      &&\text{by $\GF(2)$-linearity}\\
    &= L(r_0 \oplus (r_0 \wedge \neg m)) \oplus w \\
    &= L(r_0 \wedge m) \oplus w
      &&\text{since $r_0 \oplus (r_0 \wedge \neg m)
      = r_0 \wedge m$} \\
    &= L(r_0 \wedge m) \oplus (r_0 \wedge \neg m).
\end{align*}

The identity $r_0 \oplus (r_0 \wedge \neg m) = r_0
\wedge m$ holds bitwise: for a single bit $b$ and
mask bit $n$, $b \oplus (b \wedge (1 \oplus n)) =
b \wedge n$, verified by the truth table.

\emph{Interpretation.} $L(r_0 \wedge m)$ is $L$
applied only to the masked lanes, and $r_0 \wedge
\neg m$ is the original register's value on the
unmasked lanes. The XOR combines them lane-by-lane
with no interference, since mask-disjoint lanes
contribute nonzero bits to only one term.

%----------------------------------------------------------------------
\subsection{Fused-apply correctness}\label{app:b-fused}
%----------------------------------------------------------------------

Full verification of Proposition
\ref{prop:fused-correctness}.

\emph{Setup.} Atoms $A_1, \ldots, A_n$, each
$\GF(2)$-linear on its assigned lane range
$\mathrm{range}(A_i) \subseteq [N]$. Masks $m_i$ with
$\mathrm{supp}(m_i) = \mathrm{range}(A_i)$, pairwise
disjoint. Fused matrix $M =
\mathrm{blockdiag}(M_{A_1}, \ldots, M_{A_n})$.

\emph{Claim.} For any input $r \in \GF(2)^N$,
$(M \cdot r)|_{\mathrm{range}(A_i)} =
A_i(r|_{\mathrm{range}(A_i)})$.

\emph{Proof.} Let $\ell \in \mathrm{range}(A_i)$.
By block-diagonal structure, $M[\ell, k] = 0$ for
$k \notin \mathrm{range}(A_i)$. Hence
\begin{equation}
  (M r)[\ell] = \bigoplus_k M[\ell, k] \cdot r[k]
  = \bigoplus_{k \in \mathrm{range}(A_i)}
    M_{A_i}[\ell, k] \cdot r[k]
  = (M_{A_i} \cdot r|_{\mathrm{range}(A_i)})[\ell]
  = A_i(r|_{\mathrm{range}(A_i)})[\ell].
\end{equation}

For $\ell$ outside every range, $M[\ell, :] = 0$
(no block covers this row), so $(M r)[\ell] = 0$,
the identity under $\oplus$. The unused lanes carry
zeros, consistent with ``no atom acts here.''

%----------------------------------------------------------------------
\subsection{Cost-bounded Earley completeness}\label{app:b-earley}
%----------------------------------------------------------------------

We verify Theorem \ref{thm:earley-complete} by
reduction to A$^*$ completeness.

\emph{Setup.} The parse forest of $(W, G)$ is viewed
as a directed graph whose nodes are chart items and
whose edges are one-step production applications.
Start node: the empty item
$(S \to \bullet \delta, 0, 0, \bar{a}_0, \bar{0})$
for the grammar's start symbol $S$. Goal nodes:
items $(S \to \delta \bullet, 0, N, \bar{a},
\bar{c})$ with $\bar{c} \leq B$ componentwise.

\emph{Costs.} Each edge's cost is the incremental
cost vector added by the production application.

\emph{Heuristic.} The admissible $h^*$ of
Proposition \ref{prop:heuristic-monotone} lower
bounds the cost-to-goal from any chart item.

\emph{A$^*$ correctness argument.} With an
admissible, monotone heuristic, A$^*$ on the parse
forest graph:
\begin{itemize}[nosep]
  \item never expands a node whose $f = g + h$
    exceeds the cost of an actual goal path;
  \item terminates (the graph is finite: chart cells
    are bounded by $N^2$, productions are finite,
    attribute tuples are bounded by dominance
    pruning);
  \item returns an optimal goal path.
\end{itemize}

The cost-bounded completer's rejection is exactly
A$^*$'s pruning of nodes with $f > B$: it is sound
because any rejected item could only extend to
plans with cost $\geq$ its current cost (by
monotonicity of the cost combinator), and such
plans already violate the budget.

\emph{Completeness.} If any admissible plan exists,
A$^*$ finds one. If none exists, A$^*$ terminates
with the chart empty at $\chi[0, N]$'s goal cell,
and the planner reports infeasibility.

%----------------------------------------------------------------------
\subsection{Lagrangian achievability, affine case}
\label{app:b-lagrangian}
%----------------------------------------------------------------------

We prove the affine-cost special case of Theorem
\ref{thm:lagrangian}.

\emph{Setup.} Each production's cost is an affine
function of batch size $k$: $\bar{c}(k) = \bar{c}_0
+ k \bar{c}_1$. This holds for base productions
(atoms scale linearly in throughput-dominated
resources and sublinearly in latency-dominated
ones, but the subgradient update treats both as
affine).

\emph{Claim.} The LP relaxation of the primal
has a unique optimum, and its dual has a unique
optimum; the primal plan constructed from the
dual-optimal price is LP-optimal.

\emph{Proof.} The LP is
\begin{equation}
  \min \bar{c}_{\LAT} \cdot x
  \quad \text{s.t.} \quad
  A x \leq B, \quad x \geq 0,
\end{equation}
where $x$ is an indicator vector over
derivation paths, $A$ is the constraint matrix
extracting the resource-dimensional cost of a plan,
and the objective is the latency component.

The LP has standard strong duality: primal
optimal equals dual optimal. The dual-optimal
$\lambda^*$ is a non-negative vector with
complementary slackness: $\lambda_r^* \cdot
(A x - B)_r = 0$. Running the weighted parse at
$\lambda^*$ finds a $P$ that minimizes
$\bar{c}_{\LAT} + \lambda^* \cdot \bar{c}_{\neg
\LAT}$, which is the Lagrangian at the
dual-optimal $\lambda$, and by strong duality
equals the LP optimum.

\emph{Integer gap.} The LP relaxation allows
fractional $x$, which corresponds to convex
combinations of plans --- a notion the execution
hardware cannot realize. The gap between the LP
optimum and the IP optimum (the best \emph{single}
plan) is bounded in the affine case by the LP
sensitivity $\partial \bar{c}_{\LAT} / \partial x$
at the optimum, and is zero when the fractional
optimum happens to lie at an integer vertex of
the feasible polytope. Experimentally for our
workloads the gap is always $\leq 2$.

%----------------------------------------------------------------------
\subsection{LCM density counting}\label{app:b-lcm}
%----------------------------------------------------------------------

Proposition \ref{prop:lcm-density} restated: in an
$L \times L$ block, with $L = \lcm(d, M)$, the CSTZ
bit count equals $L^2$, and the hardware-tile count
is $(L/M)^2$, so bits per hardware tile equal
$L^2 / (L/M)^2 = M^2$.

\emph{Numerical verification at $d = 3, M = 16$.}
$L = 48$. CSTZ tiles: $(48/3)^2 = 256$. Hardware
tiles: $(48/16)^2 = 9$. CSTZ tiles per hardware
tile: $256/9 = 28.44$ (average; integer counts
per tile range over $\{28, 29\}$ because
$256 = 9 \cdot 28 + 4$). Bits per hardware tile:
$256 \cdot 9 / 9 = 256 = M^2$. $\checkmark$

\emph{At $d = 2, M = 16$.} $L = 16$. CSTZ tiles:
$(16/2)^2 = 64$. Hardware tiles: 1. CSTZ tiles per
hardware tile: 64. Bits per hardware tile:
$64 \cdot 4 = 256 = M^2$. $\checkmark$

\emph{At $d = 5, M = 16$.} $L = 80$. CSTZ tiles:
$(80/5)^2 = 256$. Hardware tiles: $(80/16)^2 = 25$.
CSTZ tiles per hardware tile: $256/25 = 10.24$.
Bits per hardware tile: $256 \cdot 25 / 25 \cdot 25
/ 25 = 256$. $\checkmark$

The verification reduces to the arithmetic
identity $d^2 \cdot (L/d)^2 / (L/M)^2 = M^2$,
which holds for any $d, M$ with $L = \lcm(d, M)$.

%----------------------------------------------------------------------
\subsection{Register-time-space bound}\label{app:b-rts}
%----------------------------------------------------------------------

Theorem \ref{thm:reg-bound} follows from the
strip-packing theorem (Coffman, Garey, Johnson,
Tarjan 1980).

\emph{Setup.} Each plan $P_i$ in a portfolio claims
a $w_i \times t_i$ rectangle of the register-time
grid (width $=$ register lanes, height $=$
dispatch cycles).

\emph{Lower bound.} Area conservation:
$\sum_i w_i t_i \leq W T$ where $W$ is the total
register width and $T$ the total cycle count.

\emph{Upper bound.} Next-Fit-Decreasing-Height
(NFDH) strip packing produces a packing with
height $T \leq 2 \cdot T^*$ where $T^*$ is the
optimal packing height for a strip of width $W$.
This is the classical 2-approximation (Coffman
et al.\ 1980; Kenyon \& R\'{e}mila 2000 give
asymptotic-optimal algorithms if tighter bounds
are needed).

In our setting the width is fixed (the register
file), and we want the cycle count to be as small
as possible; NFDH gives the 2-approximate answer
in $O(k \log k)$ time on $k$ plans.
